\documentclass[11pt,a4paper]{article}
\usepackage[utf8]{inputenc}
\usepackage[T1]{fontenc}
\usepackage[turkish]{babel}
\usepackage[margin=2.2cm]{geometry}
\usepackage{array}
\usepackage{longtable}
\usepackage{hyperref}

\title{IK ROS2 Projesinden MPC ROS2 Yapisine Gecis Notu}
\author{}
\date{}

\begin{document}
\maketitle

\section*{Amac}
Elimizde IK ile calisan bir ROS2 proje var. Bunu MPC tabanli yapiya gecirirken en kritik degisiklik, sadece ``joint hedefi publish eden'' bir dugum yerine, artik \textbf{durum bilgisini MPC'ye veren} ve \textbf{MPC'den gelen joint komutunu robota uygulayan} bir kopru dugume ihtiyac olmasidir.

\section*{Eski IK Yapisinda Tipik Durum}
Genelde akıs su sekildedir:
\begin{itemize}
  \item Giris: hedef pozisyon / ayak hedefi / IK referansi
  \item Cikis: direkt joint position veya joint velocity komutu
  \item Robot tarafinda gereken sey: encoder okuyup basit geri besleme yapmak
\end{itemize}

\section*{MPC'ye Gecince Ne Degisiyor?}
MPC yalnizca hedef istemez; robota ait anlik durumu da ister.

Bu repo'da MPC tarafi su verileri bekliyor:
\begin{itemize}
  \item IMU quaternion
  \item IMU angular velocity
  \item IMU linear acceleration
  \item 12 eklem pozisyonu
  \item 12 eklem hizi
  \item 4 adet temas bilgisi
  \item zaman bilgisi
\end{itemize}

Yani artik sadece ``komut publish eden'' bir node yetmez. \textbf{sensor -> estimator/MPC -> actuator} zinciri gerekir.

\section*{Yeni ROS2 Giris / Cikis Yapisi}
\subsection*{1. Robottan MPC'ye giden veri}
Gercek robot bridge node'unun su topic'i publish etmesi gerekir:

\begin{center}
\texttt{/simulator\_sensor\_data}
\end{center}

Mesaj tipi:
\begin{center}
\texttt{legged\_msgs/msg/SimulatorSensorData}
\end{center}

Icerigi:
\begin{itemize}
  \item \texttt{simulation\_time}
  \item \texttt{imu\_quat\_values}
  \item \texttt{imu\_angvel\_values}
  \item \texttt{imu\_linacc\_values}
  \item \texttt{joint\_position\_values}
  \item \texttt{joint\_velocity\_values}
  \item \texttt{contact\_flags}
\end{itemize}

\subsection*{2. MPC'den robota gelen komut}
MPC tarafi su topic'i publish eder:

\begin{center}
\texttt{/joint\_control\_data}
\end{center}

Mesaj tipi:
\begin{center}
\texttt{legged\_msgs/msg/JointControlData}
\end{center}

Icerigi:
\begin{itemize}
  \item \texttt{joint\_torque[12]}
  \item \texttt{joint\_position[12]}
  \item \texttt{joint\_velocity[12]}
  \item \texttt{actuator\_mode}
\end{itemize}

Yani artik cikis sadece ``position command'' degil; sistem ayni anda torque + position + velocity referansi gonderebiliyor.

\subsection*{3. Baslangic icin servis}
Controller baslarken su servisi cagiriyor:

\begin{center}
\texttt{/start\_control}
\end{center}

Servis tipi:
\begin{center}
\texttt{legged\_msgs/srv/StartControl}
\end{center}

Bu servis ilk durum bilgisini donuyor. Gercek robotta simulator yerine sizin bridge node bunu saglamali.

\section*{IK'den MPC'ye Gecerken Sizin Tarafta Ne Degisecek?}
\begin{longtable}{|p{0.28\textwidth}|p{0.29\textwidth}|p{0.29\textwidth}|}
\hline
\textbf{Konu} & \textbf{IK Projesi} & \textbf{MPC Projesi} \\
\hline
Ana giris & hedef / referans & sensor verisi + referans \\
\hline
Ana cikis & joint pos / vel & torque + joint pos + joint vel \\
\hline
Robot feedback & opsiyonel / sinirli & zorunlu \\
\hline
IMU kullanimi & bazen yok & zorunluya yakin \\
\hline
Contact bilgisi & cogunlukla yok & estimator icin gerekli \\
\hline
Startup & node baslar calisir & servis ile ilk state verilmesi gerekir \\
\hline
\end{longtable}

\section*{Gercek Robot Icin Minimum Bridge}
Elinizde zaten IK ile calisan ROS2 proje varsa, yeni node veya mevcut hardware interface icinde su 3 gorev yeterli olur:

\begin{itemize}
  \item Encoder + IMU + contact verisini okuyup \texttt{/simulator\_sensor\_data} publish etmek
  \item \texttt{/joint\_control\_data} subscribe edip motor surucusune yazmak
  \item \texttt{/start\_control} servisine cevap vermek
\end{itemize}

\section*{Dikkat Edilecek Noktalar}
\begin{itemize}
  \item Quaternion sirasi bu kodda kritik: \texttt{[w, x, y, z]}
  \item Contact flag sirasi: \texttt{\{LF, RF, LH, RH\}}
  \item Joint dizilim sirasi model ile birebir ayni olmali
  \item \texttt{simulation\_time} artan ve tutarli olmali
  \item IMU frame'i trunk/base frame ile uyumlu olmali; degilse publish etmeden once donusturulmeli
\end{itemize}

\section*{SimulatorSensorData Icerigi ve Dogru Sirasi}
Bu mesajin icini gercek robotta birebir dogru doldurmak gerekir.

\subsection*{Alanlar}
\begin{itemize}
  \item \texttt{simulation\_time}
  \item \texttt{imu\_quat\_values}
  \item \texttt{imu\_angvel\_values}
  \item \texttt{imu\_linacc\_values}
  \item \texttt{joint\_position\_values}
  \item \texttt{joint\_velocity\_values}
  \item \texttt{contact\_flags}
\end{itemize}

\subsection*{Dogru format}
\begin{itemize}
  \item \texttt{simulation\_time}
  \begin{itemize}
    \item Tip: \texttt{float64}
    \item Birim: saniye
    \item Her mesajda artan zaman olmali
  \end{itemize}
  \item \texttt{imu\_quat\_values}
  \begin{itemize}
    \item 4 eleman
    \item Sira: \texttt{[w, x, y, z]}
  \end{itemize}
  \item \texttt{imu\_angvel\_values}
  \begin{itemize}
    \item 3 eleman
    \item Sira: \texttt{[wx, wy, wz]}
    \item Frame: IMU/body/local
    \item Birim: rad/s
  \end{itemize}
  \item \texttt{imu\_linacc\_values}
  \begin{itemize}
    \item 3 eleman
    \item Sira: \texttt{[ax, ay, az]}
    \item Frame: IMU/body/local
    \item Birim: m/s$^2$
  \end{itemize}
  \item \texttt{joint\_position\_values}
  \begin{itemize}
    \item 12 eleman
    \item Sira:
    \texttt{[LF\_HAA, LF\_HFE, LF\_KFE, LH\_HAA, LH\_HFE, LH\_KFE, RF\_HAA, RF\_HFE, RF\_KFE, RH\_HAA, RH\_HFE, RH\_KFE]}
  \end{itemize}
  \item \texttt{joint\_velocity\_values}
  \begin{itemize}
    \item 12 eleman
    \item Sirasi, \texttt{joint\_position\_values} ile tamamen ayni olmali
  \end{itemize}
  \item \texttt{contact\_flags}
  \begin{itemize}
    \item 4 eleman
    \item Sira: \texttt{[LF, RF, LH, RH]}
    \item Deger: \texttt{true = temas var}, \texttt{false = temas yok}
  \end{itemize}
\end{itemize}

\subsection*{Tek satirlik kontrol formu}
\begin{verbatim}
simulation_time: t
imu_quat_values:   [w, x, y, z]
imu_angvel_values: [wx, wy, wz]
imu_linacc_values: [ax, ay, az]

joint_position_values:
[LF_HAA, LF_HFE, LF_KFE,
 LH_HAA, LH_HFE, LH_KFE,
 RF_HAA, RF_HFE, RF_KFE,
 RH_HAA, RH_HFE, RH_KFE]

joint_velocity_values:
[LF_HAA, LF_HFE, LF_KFE,
 LH_HAA, LH_HFE, LH_KFE,
 RF_HAA, RF_HFE, RF_KFE,
 RH_HAA, RH_HFE, RH_KFE]

contact_flags:
[LF, RF, LH, RH]
\end{verbatim}

\section*{Kisa Ozet}
IK'de ana is ``hedeften joint komutu uretmek'' idi. MPC'de ana is:
\begin{enumerate}
  \item robottan state okumak,
  \item bunu uygun ROS2 mesaji ile MPC'ye vermek,
  \item MPC'nin joint komutunu alip motora uygulamak.
\end{enumerate}

Yani pratikte en buyuk degisiklik \textbf{publish tarafi degil, state feedback ve bridge katmani} oluyor.

\end{document}
